\begin{abstract}
\noindent
Large language model (LLM) agents are beginning to operate inside high-performance computing (HPC) environments: monitoring Slurm jobs, debugging failed builds, analyzing simulation output, and orchestrating multi-stage scientific workflows. These agents run with the credentials and scheduler authority of the users who launch them. That is precisely what makes them dangerous in a new way. An HPC agent does not need a stolen password or a kernel exploit to cause harm --- it only needs to read a poisoned log, a compromised tool description, or a message from another agent, and it can be redirected into actions that remain fully authenticated and authorized at the account level while violating the boundary of the task it was asked to perform. We call this the \emph{hijacked authorized agent} problem. We argue that this problem is not adequately addressed by either of the two literatures it sits between. LLM-agent security research has produced strong general threat models, benchmarks, and mechanisms for indirect prompt injection and tool misuse, but almost entirely in web, enterprise, and personal-assistant settings --- not in scheduler-managed, multi-project, shared-filesystem environments. HPC security research has produced mature controls for identity, isolation, and workflow security, but assumes a human is the source of intent behind every authenticated action --- an assumption that no longer holds once an LLM agent is interposed between the user and the cluster. This paper defines the hijacked authorized agent threat model for HPC, taxonomizes the HPC-specific attack surfaces it creates, analyzes why existing HPC controls do not close the gap, and lays out a research agenda --- including an empirical benchmark now in development --- for studying and mitigating it.
\end{abstract}
